% =============================================================================
% Project : THE_PHILOSOPHY_OF_VIABILITY
% Book    : Book I — Ontology of Viability
% Book ID : VS.B01
% Part    : Part I — Identity and Manifestation (VS.B01.P01)
% Chapter : Chapter 07 — Formal and Logical Reconstruction
% Chapter ID : VS.B01.P01.C07
%
% File    : 4_ASSETS/diagrams/VS.B01.P01.C07.S03.SS02.P001.DG01.domain_architecture_and_interpretation.tex
%
% Section / Chapter / Asset:
%   Diagram DG01 — Domain Architecture and Interpretation
%
% Structural Identifier : VS.B01.P01.C07.S03.SS02.P001.DG01
% Asset Type : DG
%
% Purpose :
%   To distinguish the domain of discourse, the formal domain of
%   interpretation, and the ontological domain without identifying their
%   status.
%
% Representational limit :
%   The arrows represent methodological and interpretive relations, not
%   ontological production or equivalence.
%
% Dependencies :
%   TikZ; libraries used: arrows.meta, positioning. All are loaded by the
%   project preamble; no per-file package or library loading.
%
% =============================================================================

\begin{figure}[htbp]
    \centering
    \hyphenpenalty=10000 \exhyphenpenalty=10000\relax  % no hyphenation inside diagram labels

    \begin{tikzpicture}[
        font=\small,
        node distance=1.4cm,
        % --- base box ---
        box/.style={
            draw, rounded corners=3pt, align=center,
            text width=8.2cm, minimum height=1.9cm,
            inner sep=7pt, line width=0.8pt
        },
        % --- node roles (names avoid clashes with built-in TikZ keys) ---
        discbox/.style={box, fill=blue!8,  draw=blue!55!black},
        formbox/.style={box, fill=blue!16, draw=blue!55!black},
        ontbox/.style={box, dashed, fill=gray!10, draw=black!60},
        % --- arrows (methodological and interpretive, not production) ---
        directed/.style={-{Stealth[length=2.5mm]}, line width=0.9pt},
        qualified/.style={{Stealth[length=2.5mm]}-{Stealth[length=2.5mm]},
                          line width=0.9pt, dashed},
        elabel/.style={font=\footnotesize\itshape, text=black!75,
                       align=left, text width=3.4cm}
    ]

    % ---------------------------------------------------------------
    % Three senses of domain (stacked; no identity of status asserted)
    % ---------------------------------------------------------------
    \node[discbox] (discourse) {
        \textbf{Domain of Discourse}\\[3pt]
        Subject matter selected for a particular inquiry\\
        Existent, identity, determination, relation, manifestation
    };
    \node[formbox, below=of discourse] (formal) {
        \textbf{Formal Domain of Interpretation}\\[3pt]
        Variable ranges, formal objects, assigned interpretations\\
        Domains, sorts, predicates, relations, functions
    };
    \node[ontbox, below=of formal] (ontological) {
        \textbf{Ontological Domain}\\[3pt]
        A proposed range or category of what exists\\
        Requires independent ontological justification
    };

    \draw[directed] (discourse) -- (formal)
        node[elabel, midway, right=5pt] {guides formal representation};
    \draw[qualified] (formal) -- (ontological)
        node[elabel, midway, right=5pt] {interpretation may refer to ontological claims};

    % ---------------------------------------------------------------
    % Legend and note
    % ---------------------------------------------------------------
    \node[below=0.6cm of ontological, font=\scriptsize] (legend) {
        \tikz[baseline=-0.5ex]\draw[directed] (0,0) -- (0.7,0);~methodological guidance
        \quad\quad
        \tikz[baseline=-0.5ex]\draw[qualified] (0,0) -- (0.7,0);~qualified interpretive relation
    };

    \node[below=0.4cm of legend, text width=10cm, align=center,
          font=\footnotesize\itshape] {
        Formal domains and sorts do not, by themselves,
        establish the divisions or completeness of reality.
    };

    \end{tikzpicture}

    \caption{Three senses of domain in formal ontological reconstruction.
    The arrows represent methodological and interpretive
    relations, not ontological production or equivalence.}
    \label{fig:ch07_domain_architecture}
\end{figure}

% =============================================================================
% END OF DIAGRAM DG01
% =============================================================================
